% Source: https://artofproblemsolving.com/community/c3h2323115_arr0ws_firing_range
% Stable identity: aops-topic-2323115-q07
From the second equation $y=\dfrac{12-3x}{4}$. Multiplying the first equation by $144$ gives $9x^2+16y^2=144$, and substituting:
\[9x^2+(12-3x)^2=144\ \Longrightarrow\ 18x^2-72x=0\ \Longrightarrow\ x=0\text{ or }x=4.\]
Then $y=3$ or $y=0$. The solutions are $\boxed{(x,y)=(0,3)\text{ and }(4,0)}$: the line passes through the ends of two axes of the ellipse.

\begin{asy}
usepackage("amsmath");
size(6cm);
defaultpen(fontsize(10pt));
import graph;
draw((-5,0)--(7.4,0),Arrow(6)); draw((0,-3.6)--(0,4.8),Arrow(6));
label("$x$",(7.4,0),E); label("$y$",(0,4.8),N);
draw(ellipse((0,0),4,3),blue+linewidth(0.9));
draw((-1.2,3.9)--(5.6,-1.2),red+linewidth(0.9));
dot((4,0)); dot((0,3));
label("$(4,0)$",(4,0),NE,fontsize(9pt)); label("$(0,3)$",(0,3),NE,fontsize(9pt));
label("$\dfrac{x^2}{16}+\dfrac{y^2}{9}=1$",(-2.9,-2.2),SW,blue+fontsize(9pt));
label("$3x+4y=12$",(5.6,-1.2),E,red+fontsize(9pt));
\end{asy}
